\begin{textsample}{POS Dim 3 – ma – Score 12.00 – ma/ma\_inf\_6.txt}  \label{ex:f3_pos_165}
3.
Traços, \textbf{sons}, cores e formas O campo de experiências “ traços, \textbf{sons}, cores e formas ” aponta para a relevância de ambientes que estimulem a criatividade das \textbf{crianças}, a exploração e a valorização da multissensorialidade, o protagonismo e o prazer contínuo das \textbf{crianças} pelas descobertas.
Na BNCC, esse campo é apresentado da seguinte forma : Conviver com diferentes manifestações artísticas, culturais e científicas, locais e universais, no cotidiano da instituição escolar, possibilita às \textbf{crianças}, por meio de experiências diversificadas, vivenciar diversas formas de expressão e linguagens, como as artes visuais ( \textbf{pintura}, modelagem, colagem, fotografia etc. ), a música, o teatro, a dança e o audiovisual, entre outras.
Com base nessas experiências, elas se expressam por várias linguagens, criando suas próprias produções artísticas ou culturais, exercitando a autoria ( coletiva e individual ) com \textbf{sons}, traços, \textbf{gestos}, danças, \textbf{mímicas}, encenações, canções, desenhos, modelagens, manipulação de diversos materiais e de recursos tecnológicos.
Essas experiências contribuem para que, desde muito \textbf{pequenas}, as \textbf{crianças} desenvolvam senso estético e crítico, o conhecimento de si mesmas, dos outros e da realidade que as cerca.
Portanto, a Educação \textbf{Infantil} precisa promover a participação das \textbf{crianças} em tempos e espaços para a produção, manifestação e apreciação artística, de modo a favorecer o desenvolvimento da sensibilidade, da criatividade e da expressão pessoal das \textbf{crianças}, permitindo que se apropriem e reconfigurem, permanentemente, a cultura e potencializem suas singularidades, ao ampliar repertórios e interpretar suas experiências e vivências artísticas ( BRASIL, 2017:39 ).
Esse campo comporta experiências com as múltiplas linguagens e suas formas de expressão, que necessitam de ambientes ricos de significados, que se constituem de imagens, cores, \textbf{sons}, traços e que compõem a diversidade de linguagens, as quais as \textbf{crianças} utilizam para se expressar, se comunicar e interagir com o meio. 4. \textbf{Escuta}, fala, pensamento e imaginação O campo de experiências “ \textbf{escuta}, fala, pensamento e imaginação ” envolve a oralidade, a \textbf{escuta}, o estímulo ao pensamento e à imaginação, que devem ser fomentados na Educação \textbf{Infantil}.
Isso ocorre, entre outras iniciativas, por meio da participação das \textbf{crianças} em diversificadas experiências com a língua materna.
Conforme exposto na BNCC : Desde o nascimento, as \textbf{crianças} participam de situações comunicativas cotidianas com as pessoas com as quais interagem.
As primeiras formas de interação do \textbf{bebê} são os movimentos do seu corpo, o olhar, a postura corporal, o sorriso, o choro e outros recursos vocais, que ganham sentido com a interpretação do outro. \textbf{Progressivamente}, as \textbf{crianças} vão ampliando e enriquecendo seu vocabulário e demais recursos de expressão e de compreensão, apropriando -se da língua materna – que se torna, pouco a pouco, seu veículo privilegiado de interação.
Na Educação \textbf{Infantil}, é importante promover experiências nas quais as \textbf{crianças} possam falar e ouvir, potencializando sua participação na cultura oral, pois é na \textbf{escuta} de histórias, na participação em conversas, nas descrições, nas narrativas elaboradas individualmente ou em grupo e nas implicações com as múltiplas linguagens que a \textbf{criança} se constitui ativamente como sujeito singular e pertencente a um grupo social ( BRASIL, 2017:40 ).
Importante ressaltar que as \textbf{crianças} necessitam do contato com indivíduos a fim de criar vínculos e constituir um canal comunicativo.
É no convívio com o outro que as \textbf{crianças} exercitam sua fala e o desenvolvimento das demais linguagens.
Tal como posto na BNCC : Desde cedo, a \textbf{criança} manifesta \textbf{curiosidade} com relação à cultura escrita : ao ouvir e acompanhar a leitura de textos, ao observar os muitos textos que circulam no contexto familiar, comunitário e escolar, ela vai construindo sua concepção de língua escrita, reconhecendo diferentes usos sociais da escrita, dos gêneros, suportes e portadores.
Na Educação \textbf{Infantil}, a imersão na cultura escrita deve partir do que as \textbf{crianças} conhecem e das \textbf{curiosidades} que deixam transparecer.
As experiências com a literatura \textbf{infantil}, propostas pelo educador, mediador entre os textos e as \textbf{crianças}, contribuem para o desenvolvimento do gosto pela leitura, do estímulo à imaginação e da ampliação do conhecimento de mundo.
Além disso, o contato com histórias, contos, fábulas, poemas, cordéis etc. propicia a familiaridade com livros, com diferentes gêneros literários, a diferenciação entre ilustrações e escrita, a aprendizagem da direção da escrita e as formas corretas de manipulação de livros.
Nesse convívio com textos escritos, as \textbf{crianças} vão construindo hipóteses sobre a escrita que se revelam, inicialmente, em rabiscos e garatujas e, à medida que vão conhecendo letras, em escritas espontâneas, não convencionais, mas já indicativas da compreensão da escrita como sistema de representação da língua ( BRASIL, 2017:40 ).
As Diretrizes de 2009 orientam que “ na \textbf{pequena} infância, a aquisição e o domínio da linguagem verbal estão vinculados à constituição do pensamento, à fruição literária, sendo também instrumento de apropriação dos demais conhecimentos ” ( BRASIL, 2009:24 ).
A prática pedagógica precisa ter uma organização de espaços, tempos e materiais que facilitem as interações, para que as \textbf{crianças} possam se expressar, imaginar, criar, comunicar, organizar pensamentos e ideias, bem como \textbf{brincar} e trabalhar em grupo e individualmente.
Esse campo de experiências se articula com os demais.
Apesar de a escrita ser uma linguagem extremamente valorizada na sociedade atual, na Educação \textbf{Infantil} as \textbf{crianças} necessitam vivenciar todas as linguagens humanas, pois estão descobrindo ainda o mundo e precisam de muitas formas de expressão.
Por outro lado, também é preciso mencionar que por muito tempo houve incompreensões em relação ao trabalho com a leitura e escrita nessa etapa, negando, muitas vezes, o acesso a essas linguagens aos \textbf{bebês} e \textbf{crianças}, como se esses sujeitos não fossem capazes de participar de situações de letramento.
É preciso que os professores dessa etapa compreendam que é direito das \textbf{crianças} terem acesso à linguagem escrita por meio de práticas sociais com sentido de leitura e escrita.

% matched lemmas: bebê, brincar, criança, curiosidade, escuta, gesto, infantil, mímico, pequeno, pintura, progressivamente, som
\end{textsample}
